\begin{explanationbox}
\textbf{\textcolor{CExplanation}{先不要记公式：观察两个极端位置}}

想象椭圆上一点 $P$ 从最右端 $A(a,0)$ 沿椭圆走到最上端 $B(0,b)$。
\begin{itemize}[leftmargin=*,itemsep=0.35em]
  \item 在 $A$：切线是 $x=a$，所以 $d=a$；两条焦半径分别为 $a+c$ 与 $a-c$，其积为 $a^2-c^2=b^2$。于是 $\sqrt{r_1r_2}\,d=b\cdot a$。
  \item 在 $B$：切线是 $y=b$，所以 $d=b$；两条焦半径均为 $a$，于是 $\sqrt{r_1r_2}\,d=a\cdot b$。
\end{itemize}
\textbf{直观感觉：}切线离中心越远（$d$ 越大），两条焦半径的几何平均数就越小；两个量此消彼长，乘积始终锁定为 $ab$。
\textbf{注意：}不能只凭两个特殊点证明任意位置，下面给出为什么中途也不变。

\medskip
\textbf{\textcolor{CExplanation}{真正的几何桥梁：切线是焦点夹角的外角平分线}}

设 $\theta=\angle F_1PF_2$。椭圆的光学性质说明：\textbf{在点 $P$，法线平分两条焦半径形成的内角}（等价地，切线与两条焦半径所在直线形成相等的锐角）。
记两焦点到切线的垂距为 $h_1,h_2$。
由直角三角形，
\[
h_1=r_1\cos\frac{\theta}{2},\qquad
h_2=r_2\cos\frac{\theta}{2}.
\]
原点 $O$ 是两焦点的中点，并且两焦点都在切线的\textbf{同侧}（椭圆的内部侧）；平行距离按中点取平均，因此
\[
d=\frac{h_1+h_2}{2}
 =\frac{r_1+r_2}{2}\cos\frac{\theta}{2}
 =a\cos\frac{\theta}{2}.
\]
这里最后一步只用了椭圆定义 $r_1+r_2=2a$。

\medskip
\textbf{\textcolor{CExplanation}{为什么另一部分恰好产生 $b$？}}

在三角形 $F_1PF_2$ 中，已知两腰 $r_1,r_2$，底边 $F_1F_2=2c$。
余弦定理给出
\[
\begin{aligned}
(2c)^2
 &=r_1^2+r_2^2-2r_1r_2\cos\theta\\
 &=(r_1+r_2)^2-2r_1r_2(1+\cos\theta)\\
 &=4a^2-4r_1r_2\cos^2\frac{\theta}{2}.
\end{aligned}
\]
于是
\[
r_1r_2\cos^2\frac{\theta}{2}=a^2-c^2=b^2.
\]
两式一结合，就得到 $\sqrt{r_1r_2}\,d=ab$。
\textbf{核心不是神奇的代数消元，而是：椭圆定义 $r_1+r_2=2a$ + 切线角平分性质 + 余弦定理。}

\medskip
\textbf{\textcolor{CExplanation}{与熟悉的圆怎么联系？}}

圆可以看成 $a=b=R$、两焦点重合在圆心（$c=0$）的极限情况。
这时每个点的 $r_1=r_2=R$，圆的切线离圆心始终是 $d=R$，所以
\[
\sqrt{r_1r_2}\,d=R\cdot R=R^2=ab.
\]
圆里“焦半径平均数”和“切线距中心”都是常数；椭圆里两者不再各自恒定，但\textbf{乘积恒定}。
也可以说：\textbf{圆的切线距圆心等于半径，推广到椭圆后变成两个变化距离之间的补偿关系。}

\medskip
\textbf{\textcolor{CExplanation}{与另一个重要结论的联系：焦点到切线垂距积为 $b^2$}}

上面的 $h_1,h_2$ 还满足
\[
\boxed{h_1h_2=b^2},\qquad
\frac{h_1}{r_1}=\frac{h_2}{r_2}=\frac{d}{a}.
\]
于是
\[
h_1h_2=\frac{r_1r_2d^2}{a^2}=b^2.
\]
这给出一个\textbf{知识地图连接点}：\emph{焦点到切线垂距乘积定值} 与本题并非两个无关的公式，而是同一角平分几何的两种写法。若已掌握前一个结论，本题可以两步推出。

\medskip
\textbf{\textcolor{CExplanation}{为什么 $ab$ 容易记？}}

$a$ 来自椭圆定义 $r_1+r_2=2a$；$b$ 来自勾股关系 $b^2=a^2-c^2$（在焦点三角形的余弦定理中自然出现）。
顺带一提，椭圆面积是 $\pi ab$，所以这里的常数也等于“椭圆面积除以 $\pi$”；这是\textbf{数值关联}，不能仅凭面积推出本定理。

\medskip
\textbf{\textcolor{CExplanation}{考试识别信号}}
\begin{itemize}[leftmargin=*,itemsep=0.3em]
 \item 题中同时出现“过椭圆动点的切线”“中心到切线距离”“两条焦半径”——优先想到 C020。
 \item 只有横坐标 $x_0$、不想求 $y_0$ 时，直接用 $r_1r_2=a^2-\dfrac{c^2x_0^2}{a^2}$。
 \item 已知 $r_1+r_2=2a$ 且出现 $\angle F_1PF_2$，考虑余弦定理和半角公式。
 \item 一见到 $d<b$ 或 $d>a$，先判\textbf{无解}，不要机械代公式。
\end{itemize}
\end{explanationbox}
